\cvsection{Expériences professionnelles}

\begin{cventries}
  \cventry
    {Stagiaire ingénieur de recherche en IA}
    {INRIA}
    {Bordeaux, France}
    {Mars 2026 -- Septembre 2026}
    {
      \begin{cvitems}
        \item {Analyse d’un besoin de recherche documentaire et développement en autonomie d’une application fonctionnelle pour aider les chercheurs à retrouver des publications scientifiques.}
        \item {Conception d’un agent IA multi-étapes avec l’API OpenAI, prompts adaptés au cas d’usage, appels d’outils et pipeline RAG avec ChromaDB.}
        \item {Automatisation de la collecte de métadonnées de publications en ligne avec des tâches cron et transmission des résultats par webhooks.}
        \item {Développement d’API documentées avec OpenAPI et d’une interface React/Next.js permettant aux chercheurs d’utiliser l’agent et de choisir les messages transmis au modèle.}
        \item {Définition des critères d’évaluation, comparaison des LLM selon le taux de réussite de la recherche et la latence, puis synthèse des résultats et de leurs limites.}
        \item {Évaluation d’OpenClaw dans un environnement Docker isolé pour étudier les risques liés aux actions d’un agent ayant accès à des outils et à des fichiers.}
        \item {Déploiement de LLM avec vLLM sur cluster HPC, conteneurisation de composants avec Docker et mise en place d’une CI/CD sur GitHub.}
      \end{cvitems}
    }

  \cventry
    {Stagiaire ingénieur de recherche en apprentissage par renforcement}
    {Université Karunya}
    {Coimbatore, Inde}
    {Juin 2025 -- Septembre 2025}
    {
      \begin{cvitems}
        \item {Développement d’un modèle d’apprentissage par renforcement pour la navigation autonome d’un robot hospitalier, puis déploiement et tests sur Raspberry Pi.}
      \end{cvitems}
    }
\end{cventries}
